\chapter{Auctions and Market Design}\label{ch:b3:auctions-and-market-design}

Some markets are designed: a government sells radio licences, a city assigns pupils
to schools, surgeons match kidney donors to patients. The designer chooses the
rules, and the rules decide who gets what and who pays how much. This chapter
proves the central result of auction theory, shows when it fails, and follows it to
the auctions that raised billions and those that raised almost nothing. It then
turns to markets without prices --- matching pupils and schools, and patients and
donors --- where the algorithms of this chapter are in daily use.

\begin{openingchart}{Revenue of the European third-generation mobile-telephone
licence auctions, 2000 and 2001, in euros per head of population.}{Klemperer (2002),
European Economic Review 46, Table 1}
\begin{tikzpicture}
\begin{axis}[omrecord, xbar, bar width=6pt, width=0.6\linewidth,
  height=0.44\linewidth, xmin=0, xmax=720, ymin=-0.6, ymax=8.6,
  ytick={0,...,8}, yticklabels from table={figdata/bachelor-3/28-auctions-and-market-design/out/spectrum.csv}{country},
  y tick label style={font=\scriptsize}, xlabel={Euros per head},
  nodes near coords, nodes near coords style={font=\scriptsize},
  every node near coord/.append style={anchor=west}]
  \addplot[fill=omDef!60, draw=omDef] table[x=euros, y=pos] {figdata/bachelor-3/28-auctions-and-market-design/out/spectrum.csv};
\end{axis}
\end{tikzpicture}
\end{openingchart}

% ledger: b3am:kl.t1
% ledger: b3am:kl
The licences sold were much alike, and analysts expected roughly equal values per
head across Western Europe. The auctions of 2000 raised 650 euros a head in the U.K. and
615 in Germany, but 170 in the Netherlands and 20 in Switzerland, where analysts had
expected 400 to 600 a week before the auction began. The later auctions of 2001 raised
less, as valuations had fallen everywhere. The differences in 2000 come from the
rules: how many bidders entered, and whether they could coordinate.

\begin{recall}
Bayesian Nash equilibrium, the first- and second-price sealed-bid auctions, the
statement of revenue equivalence and the winner's curse are the Year 2 volume's; the
envelope theorem is the Year 1 volume's.
\end{recall}

\begin{notation}[Symbols of this chapter]\label{not:b3:auctions-and-market-design:symbols}
$n$ bidders have values $v_i$ drawn independently from a distribution $F$ with density $f$
on $[0, \bar v]$; $b(v)$ is a bid function, $r$ a \omterm{def:b3:auctions-and-market-design:reserve}{reserve price} and $v_0$ the seller's own value
of the object.
\end{notation}

\section{Auction formats}\label{sec:b3:auctions-and-market-design:formats}

\begin{definition}[English auction; Dutch auction]\label{def:b3:auctions-and-market-design:formats}
\index{English auction}\index{Dutch auction}
In an \emph{English auction} the price rises until one bidder remains, who pays the
price at which the last rival dropped out. In a \emph{Dutch auction} the price falls
until the first bidder accepts it, and that bidder pays it.
\end{definition}

The \omterm{def:b3:auctions-and-market-design:formats}{Dutch auction} is strategically the same as the first-price sealed-bid auction: a
bidder's only decision is the price at which to accept, and she learns nothing
useful before making it. With independent private values, the \omterm{def:b3:auctions-and-market-design:formats}{English auction} is
equivalent to the second-price auction: staying in up to one's value is a dominant
strategy, and the winner pays the second-highest value.

\begin{omfigure}
\begin{tikzpicture}[econ, x=5cm, y=5cm]
  \draw[econ axes] (0,0) -- (1.08,0) node[right, font=\scriptsize] {$v$};
  \draw[econ axes] (0,0) -- (0,1.05) node[above, font=\scriptsize] {bid};
  \draw[guide] (0,0) -- (1,1) node[right, font=\scriptsize] {$b = v$ (second-price)};
  \draw[curve, omThm] (0,0) -- (1,0.75) node[right, font=\scriptsize] {$n = 4$};
  \draw[curve, omDef] (0,0) -- (1,0.5) node[right, font=\scriptsize] {$n = 2$};
  \node[below, font=\scriptsize] at (1,0) {1};
\end{tikzpicture}
\par\smallskip{\small Equilibrium bids in a first-price auction with values uniform on
$[0, 1]$: $b(v) = (n - 1)v/n$. Bidders shade their bids below their values, and shade less as rivals
multiply.}
\end{omfigure}

\section{Revenue equivalence}\label{sec:b3:auctions-and-market-design:equivalence}

\begin{proposition}[Revenue equivalence]\label{prop:b3:auctions-and-market-design:revenue-equivalence}
Let $n$ risk-neutral bidders have independent private values from $F$ on $[0, \bar v]$. Any auction in
which the bidder with the highest value wins, and a bidder with value 0 pays nothing in
expectation, gives each bidder the same expected payment, and the seller the same
expected revenue, $E[\text{second-highest value}]$.
\end{proposition}

\begin{proof}
In equilibrium, let $Q(v) = F(v)^{n-1}$ be the probability that a bidder of value $v$ wins and $P(v)$ her
expected payment, so her payoff is $U(v) = vQ(v) - P(v)$. She can act as if her value were any $w$, so
$U(v) = \max_w\,[vQ(w) - P(w)]$, and by the envelope theorem $U'(v) = Q(v)$. With $U(0) = 0$, $U(v) = \int_0^v Q(x)\,dx$, and
$P(v) = vQ(v) - \int_0^v Q(x)\,dx$, which depends only on $F$ and $n$, not on the format. Expected revenue is $n\,E[P(v)]$,
the same in every such auction; in the second-price auction it is the expected
second-highest value.
\end{proof}

\begin{model}{Independent private values}
\assumes Each bidder knows her own value; values are independent draws from one
distribution; bidders are risk neutral.
\predicts All standard formats raise the same expected revenue; the Dutch and
first-price auctions are equivalent, and so are the English and second-price
auctions.
\tested In field experiments selling collectible trading cards over the
Internet, the English and second-price formats raised about the same revenue, but
the \omterm{def:b3:auctions-and-market-design:formats}{Dutch auction} raised 30~\% more than the first-price auction, a violation of the
prediction. % ledger: b3am:lr
\breaks With risk-averse bidders the first-price auction raises more; with
asymmetric bidders or common values the formats differ.
\end{model}

Risk aversion is the simplest break. In a second-price auction bidding one's value
remains dominant. In a first-price auction, a risk-averse bidder shades less, trading
some profit when she wins for a higher chance of winning, and revenue rises. Exercise 11
measures the difference.

\section{Reserve prices}\label{sec:b3:auctions-and-market-design:reserve}

\begin{definition}[Optimal reserve price]\label{def:b3:auctions-and-market-design:reserve}
\index{optimal reserve price}
A \emph{reserve price} is the lowest price at which the seller will sell. The
\emph{optimal reserve price} maximises the seller's expected revenue plus her own
value of the object when it stays unsold.
\end{definition}

\begin{proposition}[The optimal reserve]\label{prop:b3:auctions-and-market-design:reserve}
In a second-price auction with independent private values, if $v - (1 - F(v))/f(v)$ is increasing, the
optimal reserve $r$ solves $r - (1 - F(r))/f(r) = v_0$; it does not depend on $n$. With values uniform on $[0, 1]$, $r = (1 + v_0)/2$.
\end{proposition}

\begin{proof}
Let $\psi(v) = v - (1 - F(v))/f(v)$. Integrating the expected payment $P(v)$ of the proof of revenue equivalence by
parts, a bidder pays $\int_r^{\bar v}\psi(v)Q(v)f(v)\,dv$ in expectation when only values above $r$ win, so the seller's payoff is
\[
n\int_r^{\bar v}\psi(v)F(v)^{n-1}f(v)\,dv + v_0 F(r)^n .
\]
Its derivative in $r$ is $-nF(r)^{n-1}f(r)(\psi(r) - v_0)$, zero where $\psi(r) = v_0$, positive below and negative above. For the uniform, $\psi(r) = 2r - 1$.
\end{proof}

\begin{omchart}{Expected revenue of a second-price auction against the reserve
price, values uniform on $[0, 1]$, seller's own value zero.}{The formula of this section's
proof}
\begin{tikzpicture}
\begin{axis}[omrecord, width=0.6\linewidth, height=0.36\linewidth, xmin=0, xmax=1,
  xtick={0,0.25,0.5,0.75,1},
  x tick label style={/pgf/number format/fixed, /pgf/number format/precision=2},
  ymin=0, ymax=0.75, xlabel={Reserve price $r$}, ylabel={Expected revenue},
  legend pos=south west, legend style={font=\footnotesize}]
  \addplot[omDef, thick] table[x=r, y=n2] {figdata/bachelor-3/28-auctions-and-market-design/out/reserve.csv};
  \addplot[omThm, thick] table[x=r, y=n5] {figdata/bachelor-3/28-auctions-and-market-design/out/reserve.csv};
  \draw[guide] (axis cs:0.5,0) -- (axis cs:0.5,0.70);
  \legend{2 bidders, 5 bidders}
\end{axis}
\end{tikzpicture}
\end{omchart}
% figdata: reserve

With two bidders, a reserve of one half raises expected revenue from $1/3$ to $5/12$, a quarter more;
with five, the gain is small, because the second-highest value rarely falls below the
reserve. A reserve is most useful when competition is thin, which is also when it
risks leaving the object unsold.

\section{Common values and the winner's curse}\label{sec:b3:auctions-and-market-design:common}

When the object has a common value that each bidder estimates with error --- the oil
in a tract, the future profits of a licence --- winning tells the winner that her
estimate was the most optimistic. A bidder who bids her unconditional estimate
loses money on average: the winner's curse of the Year 2 volume. Bidders who
understand it shade their bids by more as rivals multiply. When values are
affiliated (one bidder's high signal makes high signals of others more likely), the
\omterm{def:b3:auctions-and-market-design:formats}{English auction} lets each bidder learn from the prices at which others drop out;
the \emph{linkage principle} states that releasing information linked to the value
raises expected revenue, and in that environment the \omterm{def:b3:auctions-and-market-design:formats}{English auction} raises more
than the second-price auction, which raises more than the first-price auction.

\section{Multi-unit auctions and the record of spectrum sales}\label{sec:b3:auctions-and-market-design:spectrum}

\begin{definition}[Simultaneous ascending auction; clock auction]\label{def:b3:auctions-and-market-design:multi-unit}
\index{simultaneous ascending auction}\index{clock auction}
In a \emph{simultaneous ascending auction} several licences are sold at once, in rounds;
prices rise on each separately and none is sold until no one bids again on any. In a
\emph{clock auction} the auctioneer announces prices and bidders state the quantities
they demand, and prices rise where demand exceeds supply.
\end{definition}

When licences are complements, a bidder who wants a pair risks winning one at a
price justified only by the pair: the exposure problem. When they are substitutes,
bidders can divide them by signalling in early rounds.

\begin{record}[Spectrum auctions: the U.K. and the Netherlands, 2000]
The U.K. sold five licences in March and April 2000, with no bidder allowed more than
one and one reserved for a new entrant. Nine new entrants bid against the four
incumbents and the auction raised 39 billion euros. Four months later the
Netherlands used nearly the same design with five incumbents and five licences. The
strongest potential entrants joined incumbents instead of bidding; one weak entrant
bid, and stopped after an incumbent threatened legal action. The auction raised less
than 3 billion euros against almost 10 billion forecast. % ledger: b3am:kl
\end{record}

The same design produced opposite results because entry differed. An ascending
auction lets a strong bidder top any bid, which deters weaker ones; a sealed-bid stage
gives them a chance, and so brings them in.

\begin{record}[The U.S. incentive auction, 2016--2017]
The United States bought spectrum back from television stations and sold it to
mobile operators in one design: a descending \omterm{def:b3:auctions-and-market-design:multi-unit}{clock auction} in which broadcasters
offered to give up or move their channels, and an ascending \omterm{def:b3:auctions-and-market-design:multi-unit}{clock auction} for the
licences. It began on 29 March 2016 and closed on 13 April 2017. Fifty winning bidders paid
19.8 billion dollars for 70 MHz of licensed spectrum, with 14 MHz more for unlicensed use; more than 10
billion dollars went to 175 broadcasters, of whom 133 relinquished their licences and 30 moved to a
lower channel. % ledger: b3am:fcc
\end{record}

\section{Matching: schools and kidneys}\label{sec:b3:auctions-and-market-design:matching}

\begin{definition}[Stable matching; deferred acceptance algorithm; strategy-proofness]\label{def:b3:auctions-and-market-design:matching}
\index{stable matching}\index{deferred acceptance algorithm}\index{strategy-proofness}
A matching of pupils to schools is \emph{stable} if no pupil prefers a school that
either has a free seat or holds a pupil of lower priority there. In the \emph{deferred
acceptance algorithm}, each unassigned pupil applies to her most preferred school
that has not yet rejected her; each school holds its best applicants up to its
capacity and rejects the rest; the algorithm stops when no pupil is rejected. A
mechanism is \emph{strategy-proof} if reporting one's true preferences is a dominant
strategy.
\end{definition}

\begin{proposition}[Deferred acceptance]\label{prop:b3:auctions-and-market-design:da}
Deferred acceptance ends with a \omterm{def:b3:auctions-and-market-design:matching}{stable matching}, and each proposer weakly prefers it
to every other \omterm{def:b3:auctions-and-market-design:matching}{stable matching}.
\end{proposition}

\begin{proof}
It ends, since each pupil applies to each school at most once. Suppose pupil $i$ prefers
school $s$ to her match. Then $i$ applied to $s$ earlier and was rejected, when $s$ held pupils
of higher priority; a school's held set only improves, so at the end $s$ is full of
pupils of higher priority than $i$: no blocking pair. For optimality, call a school
\emph{possible} for $i$ if some \omterm{def:b3:auctions-and-market-design:matching}{stable matching} gives it to her, and suppose that some
pupil is rejected by a possible school; take the first such rejection, of $i$ by $s$, and a
\omterm{def:b3:auctions-and-market-design:matching}{stable matching} $M$ that gives $s$ to $i$. At that moment $s$ is full with pupils of higher priority
than $i$, and at least one of them, $j$, is not at $s$ in $M$, since $i$ takes one of its seats there. Pupil $j$
applied to $s$, so every school she ranks above $s$ has rejected her; none of them is
possible for $j$, because no possible school has rejected anyone before. Her school in $M$ is
possible for her, so she ranks it below $s$. Then $j$ prefers $s$ to her match in $M$ and has
higher priority there than $i$: $(j, s)$ blocks $M$, a contradiction. No pupil is ever rejected by
a possible school, so each ends at the best of them.
\end{proof}

\begin{proposition}[Incentives]\label{prop:b3:auctions-and-market-design:incentives}
Deferred acceptance is strategy-proof for the proposing side.
\end{proposition}
\admitted

\begin{record}[Boston's school choice, 2005]
Boston assigned pupils by a priority mechanism that accepted each round's
applications for good: a pupil who ranked a popular school first and missed it could
find her second choice already full. Records of choices showed some parents
strategising with care and others not; many pupils left unassigned could have had
one of their stated choices with a different ranking. In July 2005 the city replaced
the mechanism with deferred acceptance, under which ranking schools truthfully is
safe; it chose deferred acceptance over \omterm{def:b3:auctions-and-market-design:ttc}{top trading cycles}, which is also
strategy-proof. % ledger: b3am:aprs
\end{record}

\begin{definition}[Top trading cycles; kidney exchange]\label{def:b3:auctions-and-market-design:ttc}
\index{top trading cycles}\index{kidney exchange}
In \emph{top trading cycles} each agent owns an object and points to the owner of her
most preferred remaining object; every cycle that forms trades along itself and
leaves; the rest point again. A \emph{kidney exchange} matches patients who have a
willing but incompatible donor, so that each donor gives to another pair's patient.
\end{definition}

\begin{proposition}[Top trading cycles]\label{prop:b3:auctions-and-market-design:ttc}
\omterm{def:b3:auctions-and-market-design:ttc}{Top trading cycles} gives a Pareto-efficient allocation and is strategy-proof.
\end{proposition}
\admitted

\begin{omfigure}
\begin{tikzpicture}[econ, x=1cm, y=1cm]
  \node[draw, rounded corners, font=\scriptsize, align=center] (a) at (0,1.6) {pair 1\\patient A, donor $a$};
  \node[draw, rounded corners, font=\scriptsize, align=center] (b) at (3.4,0) {pair 2\\patient B, donor $b$};
  \node[draw, rounded corners, font=\scriptsize, align=center] (c) at (-3.4,0) {pair 3\\patient C, donor $c$};
  \draw[->, thick, omThm] (a) -- node[above right, font=\scriptsize] {$a$ gives to B} (b);
  \draw[->, thick, omThm] (b) -- node[below, font=\scriptsize] {$b$ gives to C} (c);
  \draw[->, thick, omThm] (c) -- node[above left, font=\scriptsize] {$c$ gives to A} (a);
\end{tikzpicture}
\par\smallskip{\small A three-way \omterm{def:b3:auctions-and-market-design:ttc}{kidney exchange}. No donor can give to her own
patient, but each can give to another's; the three operations take place at once, so
that no donor withdraws after her patient is served.}
\end{omfigure}

% ledger: b3am:rsu
\omterm{def:b3:auctions-and-market-design:ttc}{Kidney exchange} began with pairs of pairs and with exchanges in which a pair gives to
someone on the waiting list in return for priority for its own patient. Run as a
clearinghouse over many pairs, with the constraints that law and medicine impose,
the design is efficient and gives no pair a reason to hide information; simulations
suggest substantial gains, both in transplants and in the quality of the matches.

\begin{record}[Kidney exchange in the United States]
Agarwal and co-authors found the American market fragmented: most exchange
transplants were arranged by single hospitals, not on the national platforms. A
hospital that submits its pairs to a platform counts only part of the gain, which
goes in part to other hospitals' patients, and the platforms rewarded hospitals for
submitting patients and donors in ways that did not correct this.
Many hospitals ran exchanges below the scale at which pools match well. Removing
these failures could raise the number of exchange transplants by 30 to 63~\%.
% ledger: b3am:aaafk.kidney
The theory of the clearinghouse was right about what a thick pool achieves; it left
out the incentives of the hospitals that decide which pairs enter it.
\end{record}

\begin{measure}[Measuring a matching market's stability]
Stability can be checked from data: for each participant, list the partners she
ranks above her match and ask whether any of them would have preferred her to its
own match or had a free place. Each such pair is a blocking pair. A mechanism whose
outputs often contain blocking pairs invites participants to arrange matches outside
it. Roth's comparison of seven regional British markets for new doctors, which had
adopted different centralised mechanisms, found that those producing unstable
matches fared no better than the decentralised markets they had replaced:
stability mattered even in those more complex markets. % ledger: b3am:roth91
\end{measure}

\section{Exercises}\label{sec:b3:auctions-and-market-design:exercises}

\begin{exercise}[$\star$]\label{exo:b3:auctions-and-market-design:1}
In the opening chart, what is the ratio of the highest to the lowest revenue per head
in 2000? And between the U.K. and the Netherlands?
\end{exercise}

\begin{exercise}[$\star$]\label{exo:b3:auctions-and-market-design:2}
Three bidders have values uniform on $[0, 1]$. Find the expected revenue of a second-price
auction.
\end{exercise}

\begin{exercise}[$\star$]\label{exo:b3:auctions-and-market-design:3}
In a first-price auction with four bidders and values uniform on $[0, 1]$, what does a bidder of
value 0.6 bid?
\end{exercise}

\begin{exercise}[$\star$]\label{exo:b3:auctions-and-market-design:4}
What result of a laboratory or field experiment would contradict revenue
equivalence? Name one such result from the chapter.
\end{exercise}

\begin{exercise}[$\star\star$]\label{exo:b3:auctions-and-market-design:5}
Values are uniform on $[0, 1]$ and the seller values the object at 0.2. Find the optimal reserve.
\end{exercise}

\begin{exercise}[$\star\star$]\label{exo:b3:auctions-and-market-design:6}
Four pupils and four schools, one seat each, have invented preferences and
priorities. Run deferred acceptance with the pupils proposing and with the schools
proposing, and compare.
\end{exercise}

\begin{exercise}[$\star\star$]\label{exo:b3:auctions-and-market-design:7}
An incumbent's value of a licence is known to be higher than an entrant's. Why might
the entrant not bid in an \omterm{def:b3:auctions-and-market-design:formats}{English auction}, and how does a final sealed-bid stage
change that?
\end{exercise}

\begin{exercise}[$\star\star$]\label{exo:b3:auctions-and-market-design:8}
A tract's value is $V$; each of two bidders sees $V$ plus an error of $+10$ or $-10$ with equal chance,
independent of the other's. A bidder sees 100. What is her expected value of $V$, and what is it
given that her signal is the higher of the two (ties broken at random)?
\end{exercise}

\begin{exercise}[$\star\star$]\label{exo:b3:auctions-and-market-design:9}
What would applicants' strategies show if a mechanism were not strategy-proof?
\end{exercise}

\begin{exercise}[$\star\star\star$]\label{exo:b3:auctions-and-market-design:10}
Prove revenue equivalence with the envelope argument, and compute the expected
payment $P(v)$ for uniform values and $n$ bidders.
\end{exercise}

\begin{exercise}[$\star\star\star$]\label{exo:b3:auctions-and-market-design:11}
Four bidders with values uniform on $[0, 1]$ have utility $x^{0.5}$ of their gain $x$. Simulate the expected
revenue of the four formats, using the equilibrium first-price bid $(n - 1)v/(n - 1 + 0.5)$.
\end{exercise}

\begin{exercise}[$\star\star\star$]\label{exo:b3:auctions-and-market-design:12}
Four patient--donor pairs: donor $a$ is compatible with B and C, $b$ with A, $c$ with D, $d$ with A and C.
Each patient prefers compatible donors in the order listed in her row: A prefers $b$ then $d$;
B prefers $a$; C prefers $d$ then $a$; D prefers $c$. Run \omterm{def:b3:auctions-and-market-design:ttc}{top trading cycles}.
\end{exercise}

\section{Problem: Brenland Sells Its Radio Spectrum}\label{sec:b3:auctions-and-market-design:problem}

\begin{problem}[{Weekend problem --- Brenland sells its radio spectrum}]\label{pb:b3:auctions-and-market-design:spectrum}
Brenland sells one national mobile licence. Four operators bid; each one's value is
uniform on $[0, 400]$ million dalens, independent of the others'. Brenmoor, meanwhile, assigns
pupils to schools. % figdata: reserve, schools

\medskip\noindent\textbf{Part I --- choosing the format.}
\begin{enumerate}
\item Find the expected revenue of an \omterm{def:b3:auctions-and-market-design:formats}{English auction}.
\item Find the equilibrium bid in a first-price auction, and its expected revenue.
\item Find the expected payoff of the winner.
\item If the operators are risk averse, which format raises more?
\item If one operator is an incumbent whose value is known to be high, which format
would you fear in the other three's place, and why?
\end{enumerate}

\medskip\noindent\textbf{Part II --- the reserve.}
\begin{enumerate}[resume]
\item Find the revenue-maximising reserve when an unsold licence is worth nothing to
the government.
\item Find the expected revenue with that reserve.
\item Find the probability that the licence goes unsold.
\item Would the reserve change with eight operators?
\item If an unsold licence can be used by a public operator worth 100 million, what is the
optimal reserve?
\end{enumerate}

\medskip\noindent\textbf{Part III --- two licences.} Brenland now sells a northern and a
southern licence. One operator values each alone at 50 million and the pair at 150.
\begin{enumerate}[resume]
\item Are the licences substitutes or complements for this operator?
\item In separate \omterm{def:b3:auctions-and-market-design:multi-unit}{simultaneous ascending auctions}, prices reach 70 on each. Should it stay
in?
\item It stays in and wins the north at 70, but the south rises to 90 and it drops out. What is
its payoff?
\item Name this problem, and one design that removes it.
\item With operators that value the licences as substitutes, what could they do in the
early rounds?
\end{enumerate}

\medskip\noindent\textbf{Part IV --- Brenmoor's schools.} Ana, Ben and Dee rank North
first, Cai ranks South first. Ana: North, South, West. Ben: North, West, South. Cai: South,
North, West. Dee: North, South, West. North has one seat and priorities Cai, Dee, Ana, Ben;
South one seat, priorities Ana, Dee, Ben, Cai; West two seats, priorities Ben, Ana, Cai,
Dee.
\begin{enumerate}[resume]
\item Run deferred acceptance with the pupils proposing, round by round.
\item Check that the matching is stable.
\item Run the old mechanism, in which each round's applications are accepted for
good, with truthful rankings.
\item Find a blocking pair in that outcome.
\item Show that Ana gains under the old mechanism by ranking South first.
\item In the deferred-acceptance matching, which two pupils would both gain by
swapping schools? Why does the algorithm not let them?
\item Which mechanism would make that swap, and at what cost?
\item What evidence from Boston supports the change?
\item Why might a city prefer stability to that swap?
\item State the result: the revenue-maximising reserve and the expected revenue with
it, and the matching that deferred acceptance produces.
\end{enumerate}
\end{problem}
